\section{Discussion}
\label{sec:discussion}

\subsection{Mass scale and velocity distribution}
\label{sec:degeneracy}

The dynamical likelihood constrains the total mass only through the ratio of
the observed velocity scale to the intrinsic scale of $p_{\rm N}(\tilde u)$.
A uniform rescaling $\tilde u\to\alpha\tilde u$ of the intrinsic distribution
is exactly compensated by $M_{\rm tot}\to M_{\rm tot}/\alpha^2$.  Without
external information, the absolute mass scale and the velocity distribution
are therefore degenerate; in our model the degeneracy is broken only by the
solar anchor (Eq.~\ref{eq:solar}), the PARSEC-relative shrinkage
(Eq.~\ref{eq:shrink}), and the functional form of Eq.~(\ref{eq:good}).
The comparison between Sects.~\ref{sec:fixedshape} and \ref{sec:binresults}
shows this directly: changing the shape from the orbital-population
reference to the posterior shape moves the MLR by up to 30\% at $M_G\approx8$.

The same degeneracy acts on the metallicity sensitivity of
Sect.~\ref{sec:slope}.  Each bin has its own shape parameters, so a change
of the mass scale in one bin can be compensated by its shape independently
of the other bins.  The fitted shapes move in the corresponding direction:
the metal-poor bin, which lies below PARSEC, has $u_c=39.6$, whereas the
metal-rich bin, which lies above PARSEC, has $u_c=27.8$
(Table~\ref{tab:shape}).  Under uniform priors on the logarithmic shape
parameters, the data prefer this combination, and the PARSEC-relative
shrinkage of Eq.~(\ref{eq:shrink}) pulls each bin towards PARSEC, so that the
regularisation biases the measured sensitivity towards the PARSEC value
rather than away from it.  A genuine metallicity dependence of the intrinsic
velocity distribution, for example through the separation or eccentricity
distributions \citep{Hwang2021,Hwang2022}, would nevertheless be traded
against the MLR of each bin, in either direction.  The measured
$\dd\ln M/\dd\mh$ is therefore the metallicity sensitivity under the
assumption that the shape family of Eq.~(\ref{eq:good}) describes each bin;
it is not independent of that assumption.

\subsection{Why the orbital-population shape is disfavoured}

The posterior shapes of the near-solar bins have a lower cut-off $u_c$ and
a broader Gaussian core (smaller $B$) than the orbital-population reference.
Several effects can produce such a difference.  The reference assumes an
eccentricity distribution and a separation range that need not match the
selected sample; wide binaries in \gaia\ have a separation-dependent
eccentricity distribution \citep{Hwang2022}.  Projection effects and the
selection on the significance of $\Delta\mu$ remove low-velocity systems,
which reshapes the observed distribution relative to the intrinsic one.
Finally, unresolved companions in hierarchical triples raise $u$ beyond the
Keplerian value for the photometric mass; if their distribution in $u$
differs from the truncated normal of Eq.~(\ref{eq:out}), they distort the
fitted normal component.  The lower $f_{\rm out}$ of the metal-rich bin is consistent with a
metallicity-dependent triple or contamination fraction, but $f_{\rm out}$ is a mixture weight between two fixed shapes and should not
be read as a physical triple fraction.

\subsection{Comparison with PARSEC}

The sign of the metallicity dependence agrees with PARSEC, and at the
metallicity of most of the sample the dynamical MLR agrees with PARSEC to
within $\approx7$\% for $M_G\lesssim9$.  Two differences stand out.  First,
the dynamical metallicity sensitivity is $1.4$--$1.8$ times the PARSEC
value.  Second, the faintest stars ($M_G\approx13$, $M\approx0.1$--$0.2~\msun$)
are 10--20\% less massive than PARSEC at $\mh<0.3$.  This regime contains few
systems (93 systems contribute to the neighbourhood of $M_G=13$ in the full
sample), so the MLR there is constrained largely by the spline
regularisation and the monotonicity of brighter stars, and PARSEC colours and luminosities of fully convective
stars carry larger systematic uncertainties.  The faint-end deficit should
therefore be regarded as tentative.  A direct comparison with empirical MLRs
of nearby M dwarfs \citep{Mann2019}, which requires a transformation between
the $K_s$ and $G$ bands, is deferred to future work.

For the wide-binary tests of gravity \citep{Chae2023,Banik2024}, which
predict Newtonian velocities from photometric masses, a metallicity dependence
of the MLR that exceeds the model prediction implies that a
metallicity-independent MLR misestimates $\sqrt{M_{\rm tot}}$ by several per
cent for the metal-poor and metal-rich tails of a sample.

\subsection{Limitations}
\label{sec:limitations}

The results are conditional on the modelling choices listed below.

\begin{itemize}
\item \emph{Shape prior range.}  In the two most populous bins the posterior
of $B$, and in the $[-0.5,0)$ bin also that of $C$, lies against the lower
edge of the explored range.  The 27-node likelihood table cannot represent
shapes outside this range, so the reported MLR for these bins is the best fit
within the range rather than an unconstrained optimum, and the true
posterior may extend to smaller $B$.  Extending the range, or replacing the
table by direct quadrature of Eq.~(\ref{eq:LN}), is required before the
absolute masses of these bins can be quoted without qualification.
\item \emph{Shape interpolation.}  The normal-component likelihood at a
general shape is a mixture of node likelihoods, not Eq.~(\ref{eq:good}) at
the interpolated parameters.  Its accuracy relative to direct quadrature,
and the sensitivity of the fits to single-precision arithmetic, have not been
tested for the full bins.
\item \emph{Initialisation.}  The two runs per bin were intended to start
from different shape configurations, but the initial values were passed to
the sampler in unconstrained coordinates without the required transformation.
The runs therefore test reproducibility under different seeds rather than
robustness to widely separated starting points.
\item \emph{Mixing in the $[0.0,0.3)$ bin.}  The two runs in this bin have
minimum ESS of 213 and 137 and $\widehat R$ up to 1.026; their intervals are
provisional.
\item \emph{Stage-1 calibration.}  The Stage-1b posterior used here
contains three divergent transitions.  The Stage-1a parameters $b_G$,
$s_z$, and $\nu$ are held at point estimates, and the cross-system
covariance induced by the shared Stage-1b parameters is not propagated.
\item \emph{Modular cut.}  The dynamics cannot update the metallicity
calibration, and the Stage-2 uncertainties do not include uncertainty in
$M_G$, extinction, age, or the PARSEC model family.
\item \emph{Metallicity bins.}  Systems are assigned to bins by their
posterior median, so systems near a boundary contribute to the bin they are
assigned to while retaining their full metallicity posterior.  Each bin
constrains a full $(M_G,\mh)$ surface, which we evaluate only at the bin
midpoint; the metal-poor and metal-rich bins contain 186 and 385 systems.
\item \emph{PARSEC dependence.}  The MLR is a regularised correction to
PARSEC, anchored at the solar point.  Residuals relative to PARSEC are
therefore not independent measurements of PARSEC errors; they are shrunk
towards zero where the data are weak.
\end{itemize}
